% Part: propositional-logic
% Chapter: propositional-logic
% Section: preliminaries

\documentclass[../../../include/open-logic-section]{subfiles}

\begin{document}

\olfileid{pl}{syn}{pre}

\olsection{مقدمات}

\begin{thm}[\emph{اصل استقرا بر !!{formula}s}]
  \ollabel{thm:induction}
  اگر ویژگی‌ای چون~$P$ برای همهٔ !!{formula}sی اتمی برقرار باشد و
  چنان باشد که
  \begin{enumerate}
    \tagitem{prvNot}{برقراری آن برای $\lnot !A$ از برقراری آن برای~$!A$
      نتیجه شود؛}{}
    \tagitem{prvAnd}{برقراری آن برای $(!A \land !B)$ از برقراری آن برای $!A$
      و~$!B$ نتیجه شود؛}{}
    \tagitem{prvOr}{برقراری آن برای $(!A \lor !B)$ از برقراری آن برای $!A$
      و~$!B$ نتیجه شود؛}{}
    \tagitem{prvIf}{برقراری آن برای $(!A \lif !B)$ از برقراری آن برای $!A$
      و~$!B$ نتیجه شود؛}{}
    \tagitem{prvIff}{برقراری آن برای $(!A \liff !B)$ از برقراری آن برای $!A$
      و~$!B$ نتیجه شود؛}{}
  \end{enumerate}
  آنگاه $P$ برای همهٔ !!{formula}s برقرار است.
\end{thm}

\begin{proof}
  بگذارید $S$ مجموعهٔ همهٔ !!{formula}sی دارای ویژگی~$P$ باشد. آشکارا
  $S \subseteq \Frm[L_0]$. مجموعهٔ~$S$ همهٔ شرط‌های
  \olref[fml]{defn:formulas} را برآورده می‌کند: همهٔ !!{formula}sی اتمی را
  در بر می‌گیرد و تحت !!{operator}s بسته است.  $\Frm[L_0]$ کوچک‌ترین ردهٔ
  از این دست است، پس $\Frm[L_0] \subseteq S$. بنابراین $\Frm[L_0] = S$، و
  هر فرمول ویژگی~$P$ را دارد.
\end{proof}

\begin{prop}\ollabel{prop:balanced}
   هر !!{formula} در~$\Frm[L_0]$ \emph{متوازن} است، به این معنا که شمار
   پرانتزهای باز آن با شمار پرانتزهای بسته برابر است.
\end{prop}

\begin{prob}
اثبات کنید \olref[pl][syn][pre]{prop:balanced}
\end{prob}

\begin{prop} \ollabel{prop:noinit}
هیچ قطعهٔ آغازینِ سره‌ای از !!a{formula}، خود !!a{formula} نیست.
\end{prop}

\begin{prob}
  اثبات کنید \olref[pl][syn][pre]{prop:noinit}
\end{prob}

\begin{prop}[خوانش یکتا]
هر !!{formula}~$!A$ در $ \Frm[L_0]$ دقیقاً یک تجزیه به یکی از
صورت‌های زیر دارد:
\begin{enumerate}
\tagitem{prvFalse}{$\lfalse$.}{}

\tagitem{prvTrue}{$\ltrue$.}{}

\item $\Obj p_n$ برای یک $\Obj p_n \in  \PVar$.

\tagitem{prvNot}{$\lnot !B$ برای یک !!{formula}~$!B$.}{}

\tagitem{prvAnd}{$(!B \land !C)$ به‌ازای !!{formula}sی $!B$ و~$!C$.}{}

\tagitem{prvOr}{$(!B \lor !C)$ به‌ازای !!{formula}sی $!B$ و~$!C$.}{}

\tagitem{prvIf}{$(!B \lif !C)$ به‌ازای !!{formula}sی $!B$ و~$!C$.}{}

\tagitem{prvIff}{$(!B \liff !C)$ به‌ازای !!{formula}sی $!B$ و~$!C$.}{}
\end{enumerate}
افزون بر این، این تجزیه \emph{یکتا} است.
\end{prop}

\begin{proof}
با استقرا بر $!A$. برای نمونه، فرض کنید $!A$ دو خوانش متمایز به‌صورت
$(!B \lif !C)$ و $(!B' \lif !C')$ دارد. در این صورت $!B$ و~$!B'$ باید
یکسان باشند (وگرنه یکی از آن دو قطعهٔ آغازینِ سره‌ای از دیگری خواهد بود)؛ پس اگر دو
خوانشِ $!A$ متمایز باشند، علت باید این باشد که $!C$ و~$!C'$ دو خوانش متمایز
از یک دنبالهٔ واحدِ نمادها هستند، که بنا بر فرض استقرا محال است.
\end{proof}


\begin{defn}[جانشینی یکنواخت]
اگر $!A$ و~$!B$ از !!{formula}s باشند و $\Obj p_i$ یک !!{propositional
variable} باشد، آنگاه $\Subst{!A}{!B}{\Obj p_i}$ حاصلِ جایگزین‌کردنِ هر
رخدادِ $\Obj p_i$ با یک رخدادِ~$!B$ در~$!A$ را نشان می‌دهد؛ به همین ترتیب،
جانشینی هم‌زمانِ $\Obj p_1$، \dots،~$\Obj p_n$ با !!{formula}sی
$!B_1$، \dots،~$!B_n$ با
$\SSubst{!A}{\subst{!B_1}{\Obj p_1},\dots,\subst{!B_n}{\Obj p_n}}$ نشان داده می‌شود.
\end{defn}

\begin{prob} برای هر یک از پنج مورد زیر، که همگی !!{formula}s هستند، تعیین
 کنید که آیا آن !!{formula} را می‌توان به‌صورت جانشینیِ \( \Subst{!A}{!B}{\Obj p_i} \)
 نوشت که در آن \( !A \) عبارت است از (i) \( \Obj p_0 \)؛ (ii) \( ( \lnot \Obj p_0 \land \Obj
 p_1) \)؛ و (iii) \( ( ( \lnot \Obj p_0 \lif \Obj p_1 ) \land \Obj
 p_2 ) \). در هر مورد، جانشینی مربوط را مشخص کنید.
  \begin{enumerate}
    \item \( \Obj p_1 \)
    \item \( ( \lnot \Obj p_0 \land \Obj p_0 ) \)
    \item \( ( ( \Obj p_0 \lor \Obj p_1 ) \land \Obj p_2 )  \)
    \item \( \lnot ( ( \Obj p_0 \lif \Obj p_1 ) \land \Obj p_2 ) \)
    \item \( (( \lnot ( \Obj p_0 \lif \Obj p_1 ) \lif ( \Obj p_0 \lor \Obj p_1 )) \land \lnot ( \Obj p_0 \land \Obj p_1 )) \)
  \end{enumerate}
\end{prob}

\begin{prob}
تعریفی از نظر ریاضی دقیق برای $\Subst{!A}{!B}{p}$ با استقرا ارائه کنید.
\end{prob}

\end{document}
